\documentclass[10pt,a4paper]{amsart}
\usepackage{amsmath,amssymb,amsthm,graphicx}
% Retained minimal preamble for the Stacks Project chapter varieties.tex.
% Verbatim excerpt of the chapter preamble: Xy-pic setup (native preamble lines 14-18), the native theorem
% declarations with the mathematical macros (native preamble lines 154-217), and the project own no-op
% definitions of the reference/slogan/history display environments (native wrapper.tex lines 1-3) that the
% counted statement uses for its slogan. Copied mechanically; nothing added, removed or redefined.
% For commutative diagrams we use Xy-pic
\usepackage[all]{xy}

% We use 2cell for 2-commutative diagrams.
\xyoption{2cell}
% Theorem environments.
%
\theoremstyle{plain}
\newtheorem{theorem}[subsection]{Theorem}
\newtheorem{proposition}[subsection]{Proposition}
\newtheorem{lemma}[subsection]{Lemma}

\theoremstyle{definition}
\newtheorem{definition}[subsection]{Definition}
\newtheorem{example}[subsection]{Example}
\newtheorem{exercise}[subsection]{Exercise}
\newtheorem{situation}[subsection]{Situation}

\theoremstyle{remark}
\newtheorem{remark}[subsection]{Remark}
\newtheorem{remarks}[subsection]{Remarks}

\numberwithin{equation}{subsection}

% Macros
%
\def\lim{\mathop{\mathrm{lim}}\nolimits}
\def\colim{\mathop{\mathrm{colim}}\nolimits}
\def\Spec{\mathop{\mathrm{Spec}}}
\def\Hom{\mathop{\mathrm{Hom}}\nolimits}
\def\Ext{\mathop{\mathrm{Ext}}\nolimits}
\def\SheafHom{\mathop{\mathcal{H}\!\mathit{om}}\nolimits}
\def\SheafExt{\mathop{\mathcal{E}\!\mathit{xt}}\nolimits}
\def\Sch{\mathit{Sch}}
\def\Mor{\mathop{\mathrm{Mor}}\nolimits}
\def\Ob{\mathop{\mathrm{Ob}}\nolimits}
\def\Sh{\mathop{\mathit{Sh}}\nolimits}
\def\NL{\mathop{N\!L}\nolimits}
\def\CH{\mathop{\mathrm{CH}}\nolimits}
\def\proetale{{pro\text{-}\acute{e}tale}}
\def\etale{{\acute{e}tale}}
\def\QCoh{\mathit{QCoh}}
\def\Ker{\mathop{\mathrm{Ker}}}
\def\Im{\mathop{\mathrm{Im}}}
\def\Coker{\mathop{\mathrm{Coker}}}
\def\Coim{\mathop{\mathrm{Coim}}}

% Boxtimes
%
\DeclareMathSymbol{\boxtimes}{\mathbin}{AMSa}{"02}

%
% Macros for moduli stacks/spaces
%
\def\QCohstack{\mathcal{QC}\!\mathit{oh}}
\def\Cohstack{\mathcal{C}\!\mathit{oh}}
\def\Spacesstack{\mathcal{S}\!\mathit{paces}}
\def\Quotfunctor{\mathrm{Quot}}
\def\Hilbfunctor{\mathrm{Hilb}}
\def\Curvesstack{\mathcal{C}\!\mathit{urves}}
\def\Polarizedstack{\mathcal{P}\!\mathit{olarized}}
\def\Complexesstack{\mathcal{C}\!\mathit{omplexes}}
% \Pic is the operator that assigns to X its picard group, usage \Pic(X)
% \Picardstack_{X/B} denotes the Picard stack of X over B
% \Picardfunctor_{X/B} denotes the Picard functor of X over B
\def\Pic{\mathop{\mathrm{Pic}}\nolimits}
\def\Picardstack{\mathcal{P}\!\mathit{ic}}
\def\Picardfunctor{\mathrm{Pic}}
\def\Deformationcategory{\mathcal{D}\!\mathit{ef}}
\newenvironment{reference}{}{}
\newenvironment{slogan}{}{}
\newenvironment{history}{}{}
\usepackage{cleveref}
\crefname{originalreference}{original reference}{original references}
\Crefname{originalreference}{Original reference}{Original references}
\makeatletter
\newlabel{divisors-lemma-characterize-OD}{{01X0}{source}{}{sourcepointer-1}{}}
\newlabel{divisors-lemma-characterize-OD@cref}{{[originalreference][0][]01X0}{[1][0][]source}{}{}{}}
\makeatother
\crefname{originallemma}{original lemma}{original lemmas}
\Crefname{originallemma}{Original lemma}{Original lemmas}
\makeatletter
\newlabel{lemma-degree-effective-Cartier-divisor}{{0AYY}{source}{}{sourcepointer-2}{}}
\newlabel{lemma-degree-effective-Cartier-divisor@cref}{{[originallemma][0][]0AYY}{[1][0][]source}{}{}{}}
\makeatother
\crefname{originallemma}{original lemma}{original lemmas}
\Crefname{originallemma}{Original lemma}{Original lemmas}
\makeatletter
\newlabel{lemma-degree-tensor-product}{{0AYX}{source}{}{sourcepointer-3}{}}
\newlabel{lemma-degree-tensor-product@cref}{{[originallemma][0][]0AYX}{[1][0][]source}{}{}{}}
\makeatother
\title{Signs and triviality criteria for the degree of an invertible sheaf on a proper curve}
\author{The Stacks Project Authors}
\date{stacks-project:varieties.tex:local-snapshot-20261011; GNU FDL 1.2 or later}
\begin{document}
\maketitle
\noindent\textit{Editorial note.} All delivered mathematical text is the author's own, copied verbatim from the pinned Stacks Project chapter file varieties.tex; no statement or proof step was rewritten, reordered or removed.\par\medskip
\noindent\textit{Editorial note.} The counted span cites three cross-chapter labels, each recorded in assembly.external\_references with its canonical tag as a navigation pointer.\par\medskip
\noindent\textit{Editorial note.} The proof cites no bibliography item inside the counted span.\par\medskip
\noindent\textit{Allowed standard prerequisites (editorial).}\begin{itemize}
\item A section of an invertible sheaf on a curve is regular; characterization of O\_X(D). Stacks Project, Divisors, Lemma divisors-lemma-characterize-OD (tag 0B3D). Realizes the invertible sheaf with a nonzero section as the sheaf associated to an effective Cartier divisor.
\item Degree of an effective Cartier divisor and multiplicativity of the degree. Stacks Project, Varieties, Lemma lemma-degree-effective-Cartier-divisor (tag 0B3E) and Lemma lemma-degree-tensor-product (tag 0AYX). The first gives the sign and vanishing criterion and the second computes the degree of the twist used in the last part.
\end{itemize}\medskip
\ifdefined\hypertarget\hypertarget{sourcepointer-1}{}\fi
\noindent\textit{Source reference.} Original reference 01X0: sources/stacks-chow/tags.txt, source lines 2501--2501. Cross-chapter standard prerequisite cited by the author inside the counted span; it is used as a pointer and its own proof is not part of this item.\par\medskip
\ifdefined\hypertarget\hypertarget{sourcepointer-2}{}\fi
\noindent\textit{Source reference.} Original lemma 0AYY: sources/stacks-chow/tags.txt, source lines 13806--13806. Cross-chapter standard prerequisite cited by the author inside the counted span; it is used as a pointer and its own proof is not part of this item.\par\medskip
\ifdefined\hypertarget\hypertarget{sourcepointer-3}{}\fi
\noindent\textit{Source reference.} Original lemma 0AYX: sources/stacks-chow/tags.txt, source lines 13805--13805. Cross-chapter standard prerequisite cited by the author inside the counted span; it is used as a pointer and its own proof is not part of this item.\par\medskip
\begin{lemma}
\label{lemma-check-invertible-sheaf-trivial}
Let $k$ be a field. Let $X$ be a proper curve over $k$.
Let $\mathcal{L}$ be an invertible $\mathcal{O}_X$-module.
\begin{enumerate}
\item If $\mathcal{L}$ has a nonzero section, then
$\deg(\mathcal{L}) \geq 0$.
\item If $\mathcal{L}$ has a nonzero section $s$ which vanishes
at a point, then $\deg(\mathcal{L}) > 0$.
\item If $\mathcal{L}$ and $\mathcal{L}^{-1}$ have nonzero sections, then
$\mathcal{L} \cong \mathcal{O}_X$.
\item If $\deg(\mathcal{L}) \leq 0$ and $\mathcal{L}$ has a nonzero
section, then $\mathcal{L} \cong \mathcal{O}_X$.
\item If $\mathcal{N} \to \mathcal{L}$ is a nonzero map of invertible
$\mathcal{O}_X$-modules, then $\deg(\mathcal{L}) \geq \deg(\mathcal{N})$
and if equality holds then it is an isomorphism.
\end{enumerate}
\end{lemma}

\begin{proof}
Let $s$ be a nonzero section of $\mathcal{L}$. Since $X$ is a curve, we
see that $s$ is a regular section. Hence there is an effective
Cartier divisor $D \subset X$ and an isomorphism
$\mathcal{L} \to \mathcal{O}_X(D)$ mapping $s$ the canonical
section $1$ of $\mathcal{O}_X(D)$, see
Divisors, Lemma \ref{divisors-lemma-characterize-OD}.
Then $\deg(\mathcal{L}) = \deg(D)$ by
Lemma \ref{lemma-degree-effective-Cartier-divisor}.
As $\deg(D) \geq 0$ and $= 0$
if and only if $D = \emptyset$, this proves (1) and (2).
In case (3) we see that $\deg(\mathcal{L}) = 0$ and
$D = \emptyset$. Similarly for (4). To see (5) apply
(1) and (4) to the invertible sheaf
$$
\mathcal{L} \otimes_{\mathcal{O}_X} \mathcal{N}^{\otimes -1} =
\SheafHom_{\mathcal{O}_X}(\mathcal{N}, \mathcal{L})
$$
which has degree $\deg(\mathcal{L}) - \deg(\mathcal{N})$
by Lemma \ref{lemma-degree-tensor-product}.
\end{proof}

\end{document}
